% =========================================================
% Remark: Closure of Open Ball Need Not Equal Closed Ball
% Source: volume-iii/analysis/metric-spaces/notes/balls/
%
% Place: after the Closed Ball definition (def:closed-ball)
% =========================================================

\subsubsection{A Subtlety: Closure of $B_r(x)$ vs.\ $\overline{B}_r(x)$}
\label{sec:closure-vs-closed-ball}

In $\mathbb{R}$ with the standard metric you know that
$\overline{(a - r, a + r)} = [a - r, a + r]$,
i.e.\ the closure of an open ball is the closed ball of the same radius.
This is so familiar that it is tempting to think it is always true.

\textbf{It is not.}

\begin{remark*}[The closure of an open ball need not equal the closed ball]
\label{rem:closure-vs-closed-ball}
In a general metric space, $\overline{B_r(x)}$ (the topological closure of the open ball)
and $\overline{B}_r(x)$ (the closed ball, i.e.\ $\{y : d(x,y) \le r\}$) can differ.
We always have
\[
\overline{B_r(x)} \;\subseteq\; \overline{B}_r(x),
\]
but the reverse inclusion can fail.
\end{remark*}

\begin{example}[Discrete metric: closure of $B_1(x)$ is just $\{x\}$]
\label{ex:closure-vs-closed-ball-discrete}
Let $X$ be any set with at least two points, and take the discrete metric:
\[
d(x,y) = \begin{cases} 0 & x = y, \\ 1 & x \neq y. \end{cases}
\]
Fix any $x \in X$ and take $r = 1$.

\medskip
\noindent
\textbf{The open ball:}
$B_1(x) = \{y \in X : d(x,y) < 1\} = \{x\}$ (only $x$ itself has $d(x,x) = 0 < 1$).

\medskip
\noindent
\textbf{The closed ball:}
$\overline{B}_1(x) = \{y \in X : d(x,y) \le 1\} = X$ (every point satisfies $d(x,y) \le 1$).

\medskip
\noindent
\textbf{The closure of the open ball:}
In the discrete metric, every singleton $\{x\}$ is open (it equals $B_{1/2}(x)$).
Therefore $\{x\}$ is already closed.
The closure of a closed set is itself:
\[
\overline{B_1(x)} = \overline{\{x\}} = \{x\}.
\]

\medskip
\noindent
\textbf{Conclusion:}
\[
\overline{B_1(x)} = \{x\}
\qquad\text{but}\qquad
\overline{B}_1(x) = X.
\]
These are wildly different whenever $|X| > 1$.
\end{example}

\begin{remark*}[Why this happens]
In $\mathbb{R}$, the closure of $B_r(x)$ picks up the boundary points
$\{x-r, x+r\}$ because sequences \emph{approaching} $x \pm r$ from inside
the ball exist: $y_n = x + r - \tfrac{1}{n} \to x + r$.

In the discrete metric, there are no sequences in $\{x\}$
that converge to any $y \neq x$: a sequence in $\{x\}$ is constant
and stays at $x$.
So the closure of $\{x\}$ gains no new points.

\medskip
\noindent
The lesson: the closure of the open ball $B_r(x)$ is determined by
\emph{which points can be approached by sequences from inside $B_r(x)$},
not by the algebraic formula $d(x,y) \le r$.
Those two notions agree in $\mathbb{R}^n$ but can diverge in abstract spaces.
\end{remark*}

\begin{remark*}[Notation caution]
Because $\overline{B_r(x)}$ and $\overline{B}_r(x)$ can differ,
the notations are intentionally distinct:
\begin{itemize}
  \item $B_r(x)$ = open ball (strict inequality $d(x,y) < r$).
  \item $\overline{B}_r(x)$ = closed ball (non-strict $d(x,y) \le r$).
  \item $\overline{B_r(x)}$ = topological closure of the open ball.
\end{itemize}
In $\mathbb{R}^n$ these last two agree, so the distinction rarely matters
there. In abstract metric spaces, keep them separate.
\end{remark*}
